\section{Propriedades estruturais dos números naturais}
A construção dos números naturais apresentada implica as seguintes propriedades estruturais.

\begin{proposition}
    Para todo $n \in \omega$, $n\subseteq \omega$.
\end{proposition}

\begin{proof}
    Procede-se por indução em $n$.
    Temos que $n\subseteq \omega$ para $n=0$, pois $0=\emptyset$.
    Suponha que $n\in \omega$ e que $n\subseteq \omega$.
    Então $S(n)=n\cup\{n\}\subseteq \omega$, o que completa a indução.
\end{proof}

\begin{corollary}
    Para quaisquer $m, n \in \omega$, $m\leq n$ se, e somente se, $m\subseteq n$.
\end{corollary}
\begin{proof}
    Para a direção da esquerda para a direita, suponha que $m\leq n$.
    Se $m=n$, então $m\subseteq n$.

    Se $m<n$, então $m\in n$.
    Vejamos que $m\subseteq n$.
    Tome $k\in m$.
    Pela proposição anterior, $k\in \omega$.
    Assim, $k<m$ e $m<n$.
    Pela propriedade transitiva, $k<n$, e, portanto, $k\in n$.
    Assim, $m\subseteq n$.

    Para a recíproca, procedemos pela contrapositiva.
    Se $m\not\leq n$, segue da tricotomia que $n<m$.
    Da direção já provada, temos que $n\subseteq m$.
    Como $m\neq n$, não pode valer $m\subseteq n$, e, portanto, $m\not\subseteq n$.
\end{proof}

Afirmações como $n \in m$, apesar de serem úteis em vários aspectos da Teoria dos Conjuntos, podem ser pouco intuitivas e soarem até mesmo como indesejáveis.
Números costumam ser pensados como ``elementos'', e não como conjuntos, de modo que seria desejável que houvesse uma lista de propriedades que caracterizassem os números naturais de modo que não dependa da construção específica adotada.

Os Axiomas de Peano \cite{peano1889arithmetices}, formulados por Giuseppe Peano em 1889 \cite{peano1889arithmetic}, constituem um sistema de axiomas muito conhecido que tem como objetivo formalizar os números naturais.

Em linguagem moderna, eles dizem que $\mathbb N$ é um conjunto munido de uma função $S$ definida em $\mathbb N$, que, junto de um objeto $0$, o seguinte é satisfeito.
\begin{enumerate}[label=(\arabic*)]
    \item $0\in \mathbb N$;
    \item para todo $n \in \mathbb N$, $S(n)\in \mathbb N$;
    \item para todo $n \in \mathbb N$, $S(n)\neq 0$;
    \item para todo $n, m \in \mathbb N$, se $S(n) = S(m)$, então $n = m$.
    \item se $A\subseteq \mathbb N$ é tal que $0\in A$ e para todo $n \in A$, $S(n)\in A$, então $A=\mathbb N$.
\end{enumerate}

A nossa construção satisfaz os axiomas (1) e (2) pelo fato de $\omega$ ser indutivo, e satisfaz o axioma (3) pelo fato de $0$ ser o conjunto vazio, e verificamos também o axioma (5).
Abaixo, verificamos o axioma (4).

\begin{proposition}
    Para quaisquer $n, m \in \omega$, se $S(n) = S(m)$, então $n = m$.
\end{proposition}
\begin{proof}
    Procederemos pela contrapositiva.
    Se $n\neq m$, então, sem perda de generalidade, $n<m$.
    Assim, $S(n)\leq m<S(m)$, e, portanto, $S(n)\neq S(m)$.
\end{proof}

Com isso, a construção de $\omega$ apresentada satisfaz os Axiomas de Peano, e, portanto, pode ser vista como uma legítima construção dos números naturais dentro da Teoria dos Conjuntos.

Porém, ressaltamos que a noção elementar de número precede a própria Teoria dos Conjuntos: conforme discutido no Capítulo~1, precisamos dela para descrever a própria linguagem de primeira ordem.
Assim, afirmações como $1\in 2$ não devem necessariamente ser vistas como uma afirmação ontológica sobre a natureza última dos números.
Uma postura mais defensável é que tal construção é uma internalização da noção de número natural interna à Teoria dos Conjuntos que captura propriedades fundamentais da noção primitiva de número.